\documentclass[compress,12pt]{beamer} %Presentation mode
% \documentclass[compress,12pt,handout]{beamer} %Handout mode

% Instead of \usetheme{ucl} use following line now given the subfolder structure
%** Look at http://tex.stackexchange.com/questions/123700/how-to-use-local-beamer-theme-in-project-source-folder
% Before: \usetheme{ucl}; Now: \usepackage{UCL-UCL-Beamer/beamerthemeucl}; analog others in this document
\usepackage{UCL-Beamer/beamerthemeucl}

% Define user colors using the RGB model
\definecolor{dunkelgrau}{rgb}{0.4,0.4,0.4}
\definecolor{grau}{rgb}{0.6,0.6,0.6}
\definecolor{hellgrau}{rgb}{0.95,0.95,0.95}

% \definecolor{headlinecolor}{RGB}{130,0,0}
% \colorlet{headlinecolor}{pink}
\colorlet{headlinecolor}{brightblue}

\definecolor{orange}{RGB}{233,77,16}
\definecolor{myBlue}{RGB}{56,32,196}
\definecolor{myGreen}{RGB}{33,222,60}

\usepackage[latin1]{inputenc}
\usepackage[english]{babel}

\usepackage
    [separate-uncertainty=true,
    %multi-part-units=single,
    list-units=repeat,per-mode = symbol]
    {siunitx}
%%% Increase the height of the banner: the argument is a scale factor >=1.0
% \setbeamertemplate{banner}[ucl][1.5]

%%% Change the colour of the main banner
%%% The background should be one of the UCL colours (except pink or white):
%%%   black,darkpurple,darkred,darkblue,darkgreen,darkbrown,richred,midred,
%%%   navyblue,midgreen,darkgrey,orange,brightblue,brightgreen,lightgrey,
%%%   lightpurple,yellow,lightblue,lightgreen,stone
% \setbeamercolor{banner}{bg=darkpurple}
% \setbeamercolor{banner}{bg=brightblue}
% \setbeamercolor{banner}{bg=yellow,fg=white}

%%% Set a background colour
% \setbeamercolor{background canvas}{bg=darkgrey}

%%% Change colour of text if desired
% \setbeamercolor{normal text}{fg=white}\usebeamercolor*{normal text}

%%% Add a stripe behind the banner
% \setbeamercolor{banner stripe}{bg=pink,fg=white}
% \setbeamercolor{banner stripe}{bg=stone,fg=white}
\setbeamercolor{banner stripe}{bg=pink!30,fg=pink}

%%% The main structural elements
% \setbeamercolor{structure}{fg=white}


%%% Author/Title/Date and slide number in the footline
% \setbeamertemplate{footline}[author title date]
% \setbeamertemplate{footline}[author title frame]
\setbeamertemplate{footline}[title author frame]
% \setbeamertemplate{slide counter}[frametotal]

%%% Puts the section/subsection in the headline
% \setbeamertemplate{headline}[section]

%%% Puts a navigation bar on top of the banner
%%% For this to work correctly, the each \section command needs to be
%%% followed by a \subsection. Requires one extra compile.
% \setbeamertemplate{headline}[miniframes]
%%% Accepts an optional argument determining the width
% \setbeamertemplate{headline}[miniframes][\paperwidth]


%%% Puts the frame title in the banner
%%% Won't work correctly with the above headline templates
% Instead of \useoutertheme{ucltitlebanner} use following line now given the subfolder structure
% \usepackage{UCL-Beamer/beamerouterthemeucltitlebanner}
%%% Similar to above, but smaller (and puts subtitle on same line as title)
% \usepackage[small]{UCL-Beamer/beamerouterthemeucltitlebanner}

\usepackage[title+sectiononlyheadline]{UCL-Beamer/beamerouterthemeucltitlebanner}
% \usepackage[title+headline]{UCL-Beamer/beamerouterthemeucltitlebanner}

%%% Gives block elements (theorems, examples) a round, shaded border
\usepackage{UCL-Beamer/beamerinnerthemeblockrounded}

%%% Gives block elements (theorems, examples) a round, but non-shaded, border
% \usepackage{UCL-Beamer/beamerinnerthemeblockroundedwithoutshadow}

%%% Gives block elements (theorems, examples) a border
% \usepackage{UCL-Beamer/beamerinnerthemeblockborder}
%%% Sets the body of block elements to be clear
% \setbeamercolor{block body}{bg=white,fg=black}

%%% Include logo on title slide
\titlegraphic{%
  \includegraphics[height=0.08\paperwidth]{UCL-Beamer/logos/EPSRC}%
  \hspace{0.5cm}%
  \includegraphics[height=0.08\paperwidth]{UCL-Beamer/logos/GIFTSurg}%
  \hspace{0.5cm}%
  \includegraphics[height=0.08\paperwidth]{UCL-Beamer/logos/EPSRC_CDT}%
  }


%%% Set a background image
%%% Some sample images are available from the UCL image store:
%%%   https://www.imagestore.ucl.ac.uk/home/start
% \setbeamertemplate{background canvas}{%
%   \includegraphics[width=\paperwidth]{imagename}}



%%%%%% Some other settings that can make things look nicer
%%% Set a smaller indent for description environment
% \setbeamersize{description width=2em}

% Remove navigation symbols
% \setbeamertemplate{navigation symbols}{}


%%% Create headline slide without framenumber
%%% Create headline slide without framenumber
\def\headline#1{
  \setbeamertemplate{footline}[title author]
  \begin{frame}[c]\centering\color{headlinecolor}\LARGE\bfseries#1\end{frame}
  \setbeamertemplate{footline}[title author frame]
  \addtocounter{framenumber}{-1} %could also use noframenumbering above but it does not show the section header in the first "headline" slide for whatever reason. This solution works.
}
\newcommand{\subheadline}[2]{
  \begin{frame}\color{headlinecolor}{#1}\begin{center}\color{headlinecolor}\Large\bfseries#2\end{center}\end{frame}
}
\newcommand{\remember}[2][0.5]{{\scriptsize\hfill$\displaystyle #2$}\vspace*{#1 cm}}
% \newcommand{\remember}[2][0.5]{{\scriptsize\hfill\color{dunkelgrau}$\displaystyle #2$}\vspace*{#1 cm}}


\title[Short Title]{Title of the Presentation}
\author[Auth. 1 \and Auth.2]{Author 1 \and Author 2}
\institute[UCL]{%
  Translational Imaging Group, CMIC\\%
  University College London
}
\date{\today}


\def\R{\mathbb{R}}
\def\N{\mathbb{N}}
\def\Z{\mathbb{Z}}
\def\P{\mathbb{P}}

\def\G{\mathbb{G}}

\def\e{\mathbf{e}}
\def\g{\mathbf{g}}
\def\h{\mathbf{h}}
\def\u{\mathbf{u}}
\def\x{\mathbf{x}}
\def\z{\mathbf{z}}
\def\c{\mathbf{c}}
\def\b{\mathbf{b}}
\def\y{\mathbf{y}}
\def\r{\mathbf{r}}
\def\f{\mathbf{f}}
\def\b{\mathbf{b}}
\def\n{\mathbf{n}}
\def\p{\mathbf{p}}
\def\0{\mathbf{0}}

\def\im{\mathrm{i}}

\def\BB{\mathbf{B}}
\def\DD{\mathbf{D}}
\def\SS{\mathbf{S}}
\def\KK{\mathbf{K}}
\def\MM{\mathbf{M}}
\def\HH{\mathbf{H}}
\def\WW{\mathbf{W}}

\def\Id{\mathbf{Id}}


%\def\u{\boldsymbol{u}}
%\def\c{\mathbf{c}}
%\def\0{\mathbf{0}}

\def\II{\mathcal{I}}

\def\E{\mathcal{E}}
\def\O{\mathcal{O}}
\def\T{\mathcal{T}}

\def\I{\texttt{I}}

\def\subset{\subseteq}
\def\squareShadow{\ding{113}}	
\def\MyLineWidth{1pt}


\newcommand\spana{\text{\rm span}}
\newcommand\argmin{\operatornamewithlimits{argmin}}
\newcommand\argmax{\operatornamewithlimits{argmax}}
\newcommand{\diam}[1]{\text{\rm diam}(#1)}
\renewcommand{\Re}{\text{\rm Re}}
\renewcommand{\Im}{\text{\rm Im}}
\newcommand{\length}{\text{\rm length}}
\newcommand{\diag}{\text{\rm diag}}
\newcommand{\quadr}{\text{\rm quad}}
\newcommand{\conv}[2]{\text{\rm conv}\{#1,#2\}}
\newcommand{\set}[2]{\left\{#1\,:\,#2\right\}}
% \newcommand{\set}[2]{\left\{#1\,\middle|\,#2\right\}}
\newcommand{\norm}[2][]{\ifthenelse{\equal{#1}{}}{\left\Vert #2 \right\Vert}{\left\Vert #2\right\Vert_{#1}}}
\newcommand{\skp}[3][]{\ifthenelse{\equal{#1}{}}{\left\langle #2\,,\,#3 \right\rangle_{L^2(\R)}}{\left\langle #2\,,\,#3 \right\rangle_{#1}}}
\newcommand{\mean}[1]{\left\langle #1\right\rangle}

\newcommand{\emean}[2][]{\langle\!\;\!\!#1\langle#2\rangle\!\;\!\!#1\rangle}

\newcommand{\GA}{\text{GA}}
\newcommand{\TR}{\text{TR}}
\newcommand{\TE}{\text{TE}}

\def\myVspaceTable{\vspace*{0.5cm}}

\newcommand{\mypsfrag}[4]{\psfrag{#4}[#1][#2]{#3 \num{#4}}}

%**Scheiss Problem!?!?
% zB: ersetze 83,84,85: er findet sie aber ersetzt alle durch letzten Eintrag (85)!?
% \newcounter{ctrpsfrag}
% \newcommand{\mypsfragloop}[6]{
% 	\forLoop[#5]{#4}{#6}{ctrpsfrag}{%
% 		\psfrag{\arabic{ctrpsfrag}}[#1][#2]{#3 $\arabic{ctrpsfrag}$}
% 	}
% }

% 	\mypsfragloop{r}{r}{\scriptsize}{83}{1}{89}

%%%%%%%%%%%%%%%%%%%%%%%%%%%%%%%%%%%%%%%%
%%% Selbst-nummerierende Konstanten: %%%
%%%%%%%%%%%%%%%%%%%%%%%%%%%%%%%%%%%%%%%%
\newcommand{\newconst}[1]{%
	\refstepcounter{const}%
	C_{\theconst}\label{const:#1}%
}
\newcommand{\const}[1]{C_{\ref{const:#1}}}


%Makro f\"ur aufeinander folgende Erh\"ohung der Lemma-/Satz-/Definitions- und Algorithmus-Counter
\newcommand{\updateCtr}[1]{
	\addtocounter{helper}{1}	%Master-Ctr inkrementieren

	\FPset\dummy{\thehelper}  %Da Befehl refstepcounter auch um eines erh\"oht, w\"urde ohne diese Zeilen
	\FPsub\dummy\thehelper{1} %erster Counter-Wert gleich 2 sein.
	\FPclip\dummy\dummy

	\setcounter{#1}{\dummy}
	\refstepcounter{#1}
}




\begin{document}


\begin{frame}
  \titlepage
\end{frame}

%%% Include logo in bottom right corner of all slides
\logo{%
  \includegraphics[width=0.12\paperwidth]{UCL-Beamer/logos/GIFTSurg}}
%%% Remove nav symbols (and shift any logo down to corner)
\setbeamertemplate{navigation symbols}{\vspace{-2ex}} 

\begin{frame}[t]\frametitle{Agenda}
   \tableofcontents%[hideallsubsections]
\end{frame}


%%%%%%%%%%%%%%%%%%%%%%%%%%%%%%%%%%%%%%%%%%%%%%%%%%%%%%%%%%%%%%%%%%%%%%%%%%%%%%%%%%%%%
%%% Introduction

\section{Motivation \& Aims}
\headline{Motivation \& Aims}

% \subsection{CVDs}

\begin{frame}[t]\frametitle{Motivation}\framesubtitle{Subtitle}
    \begin{itemize}[<+->]\setlength{\itemsep}{12pt}
        \item Scanning time considerations and patient motion limits spatial resolution in MRI
        \item In case of moving subjects:
        \begin{itemize}
            \item 2D sequences with {\bf anisotropic voxel spacing}
            \item High SNR {\bf in-plane} resolution visually appealing
            \item \textbf{Often overlooked:} Effect of partial voluming in slice-select direction
        \end{itemize}
        \item Super-resolution reconstruction (SRR) approaches are capable to reduce partial voluming effects (PVE)
    \end{itemize}
\end{frame}

\begin{frame}[t,label=Aims]\frametitle{Aims}
    \begin{itemize}\setlength{\itemsep}{12pt}
        \item<1-> Comparison and evaluation of all SRR MRI techniques
        \begin{itemize}
            \item Frequency domain
            \item Image domain
            \item Wavelet domain
        \end{itemize}
        \item<2-> Evaluation of use of all approaches to correct through-plane PVEs in MRI
        \item<3-> Focus on resolution enhancement $\Rightarrow$ {\bf motion-free} scenarios
    \end{itemize}
\end{frame}

\section{Methods}
\headline{Methods}

\subsection[Frequency]{SRR in Frequency Domain} % (fold)
\label{sub:k_space}

% subsection k_space (end)
\begin{frame}[c]\frametitle{SRR in Frequency Domain}
    \includegraphics[width=1\textwidth]{figures/Gholipour2015_fig1.png}

    \visible<2->{
    \textbf{$\boldsymbol{k}$-space sampling:}
    \begin{itemize}
        \item Phase and frequency encoding directions: HF sampling
        \item Slice-select direction: No HF sampling
    \end{itemize}
    }
    \visible<3->{
    $\Rightarrow$ Fusion of $k$-space data from different stacks!
    }
\end{frame}


\subsection[Image]{SRR in Image Domain} % (fold)
\label{sub:image_domain}

\begin{frame}{SRR in Image Domain}
    \visible<1->{
        \begin{minipage}{0.65\textwidth}
            \textbf{Slice acquisition model:}
            \begin{align*}
                \y_k = \DD_k\,\BB_k\,\SS_k\,\x + \n_k
            \end{align*}
            \vspace*{-.5cm}
            {\small
            \begin{align*}
                &\y_k   &&     \text{ Low-resolution slice } k \\
                &\x     &&     \text{ High-resolution data (sought SRR)} \\
                &\DD_k  &&     \text{ Down-sampling matrix} \\
                &\BB_k  &&     \text{ Blur matrix (PSF)} \\
                &\SS_k  &&     \text{ Slice-select operator} \\
                &\n_k   &&     \text{ Noise vector }
            \end{align*}
            }
        \end{minipage}}
    \hfill
    \begin{minipage}{0.27\textwidth}
        \includegraphics[width=1.2\textwidth]{figures/SRR_ImageDomain.png}
    \end{minipage}
\end{frame}
% subsection image_domain (end)

\subsection[Wavelet]{SRR in Wavelet Domain} % (fold)
\label{sub:wavelet_domain}
\begin{frame}[t]\frametitle{SRR in Wavelet Domain}
    \begin{minipage}{0.65\textwidth}
        \textbf{Key idea:}
        \begin{itemize}
            \item Decompose MRI scan with 3D wavelet transform:
            \begin{itemize}
                \item 8 blocks
                \item in $XYZ$-coordinate system
            \end{itemize}
        \item Low- and high-frequency information in image directions
        $\left.
      \begin{array}{l}
          LLL \\
          LLH \\
          LHL \\
          \vdots \\
          HHH 
      \end{array}
    \right\}  \Rightarrow 2^{3}=8$ blocks
        \end{itemize}
        % \vspace*{0.5cm}
        \tiny{Aganj et al. (2012). A 3D wavelet fusion approach for the reconstruction of isotropic-resolution MR images from orthogonal anisotropic-resolution scans.}
    \end{minipage}
    \hfill
    \visible<2->{
    \begin{minipage}{0.27\textwidth}
        \includegraphics[width=1.2\textwidth]{figures/SRR_WaveletDomain.png}\\
    \end{minipage}
    }
\end{frame}
% subsection wavelet_domain (end)

\subsection[Evaluation]{Evaluation Methods} % (fold)
\label{sub:evaluation_methods}

\begin{frame}[t]\frametitle{Evaluation Methods}
    % \begin{minipage}{0.45\textwidth}
        \begin{exampleblock}<1->{1) Comparison of SRR Methods: Phantom Imaging}
            \begin{itemize}
                % \item American College of Radiology phantom imaging
                % \item
                \item T2-weighted FSE MRI
                \begin{itemize}
                    \item  $4\times 4$ pinhole structures \\
                        (diameters of \SI{1.1}{\mm}, \SI{1.0}{\mm} and \SI{0.9}{\mm})
                    \item In-plane resolution $\SI{0.7}{\mm}\times\SI{0.7}{\mm}$
                    \item Slice thickness: \SI{1.4}{\mm}
                \end{itemize}
                \item Isotropic HR reference: \SI{0.7}{\mm}

            \end{itemize}
        \end{exampleblock}
    % \end{minipage}\hfill

    % \begin{minipage}{0.45\textwidth}
        \begin{alertblock}<2->{2) Clinical Application: Neuroimaging}
            \begin{itemize}

                \item Clinical protocol imaging for TSC patients
                \item T2-weighted FSE MRI
                \begin{itemize}
                    \item In-plane resolution: $\SI{0.4}{\mm}\times\SI{0.6}{\mm}$
                    \item Slice thickness: \SI{2}{\mm}
                \end{itemize}
            \end{itemize}
        \end{alertblock}
    % \end{minipage} rounded=true,shadow=true
\end{frame}

\begin{frame}[t]\frametitle{Evaluation Methods}\framesubtitle{Some more methods}
    \begin{block}<2->{2) Clinical Application: Neuroimaging}
        \begin{itemize}

            \item Clinical protocol imaging for TSC patients
            \item T2-weighted FSE MRI
            \begin{itemize}
                \item In-plane resolution: $\SI{0.4}{\mm}\times\SI{0.6}{\mm}$
                \item Slice thickness: \SI{2}{\mm}
            \end{itemize}
        \end{itemize}
    \end{block}
\end{frame}

% subsection evaluation_methods (end)

\end{document}
